\section{Introduction} \label{sec:introduction}
The goal of this project is to model and solve the \textit{movers problem} using Boolean satisfiability (SAT).
The problem consists of emptying an $n$-floor building with a team of $m$ movers.
The building contains a number of items located on different floors, and all persons and items must eventually end on the ground floor, represented as floor $0$.

\noindent Time is assumed to be discrete.
At every time step, each mover may perform one of the available actions:
\[
    \texttt{ascends}(p,t), \qquad
    \texttt{descends}(p,t), \qquad
    \texttt{carries}(p,i,t),
\]
where $p$ is a person, $i$ is an item, and $t$ is a time step.
The action \texttt{ascends} moves a person up one floor, \texttt{descends} moves a person down one floor, and \texttt{carries} represents a person carrying an item down one floor while descending.
Each action takes exactly one time step.

\noindent The project was implemented as a complete SAT-based logistics planning system.
The backend encodes the problem as a conjunctive normal form (CNF) formula and solves it using the Glucose3 SAT solver through the PySAT library.
The frontend provides an interactive graphical interface where users can configure the number of floors, movers, items, and solving mode, then visualize the computed plan step by step.

\noindent The final system supports two task variants:

\begin{itemize}
    \item \textbf{Task 1: Movement Only.} Movers start at floor $0$ and must reach specified target floors within a given time horizon.
    \item \textbf{Task 2: Movement and Carrying.} Movers and items must all end on floor $0$.
    Items can only move downward when carried by a descending mover.
\end{itemize}

\noindent The main objective of the solver is not only to check whether a plan exists for a fixed number of steps, but also to find the shortest valid plan using iterative deepening over the time horizon.